\chapter{Équations et mise en équation}
\label{exercices:06}
\section*{Exercice 6.1 -- Premier degré et cas particuliers}
Résoudre $5(2x-1)=3x+16$, $\frac x3-\frac12=\frac56$, $2x+3=2x+5$ et $4(x+1)=4x+4$.

\section*{Exercice 6.2 -- Produits et carrés}
Résoudre $(2x-3)(x+4)=0$, $x^2=45$, $x^2=-4$ et $x^2=7x$.

\section*{Exercice 6.3 -- Un rectangle à trouver}
Un rectangle a un périmètre de $58\unite{cm}$. Sa longueur dépasse sa largeur de $7\unite{cm}$. Déterminer ses dimensions puis son aire.

\section*{Exercice 6.4 -- Seuil de rentabilité}
Une carte d'activité coûte $24$ euros, puis chaque séance coûte $4$ euros. Sans carte, chaque séance coûte $7$ euros. Pour combien de séances les coûts sont-ils égaux ? À partir de combien de séances entières la carte est-elle strictement avantageuse ?
